\chapter{Fair Value}\label{ch:hf:fair-value}

The same stock quotes on a dozen venues, its index future trades in Chicago and the exchange-traded fund that holds it trades in New Jersey. The
market maker's price is none of these quotes but its own estimate, revised on every message from any of them. In this chapter's simulated
market, where a follower instrument trades on two venues and a leader moves first, the best estimate built from the follower's own book forecasts
its mid five seconds ahead 2.4\% better than the mid does. Adding the leader's price improves the forecast by 3.5\% when the leader moves half a
second first, and by 12.8\% when it moves two seconds first. Small percentages, applied to every quote of every day, are what this chapter is about.

\section{Mid, weighted mid and microprice}

\begin{definition}[Fair price]\label{def:hf:fair-value:fair}
A market maker's \emph{fair price}\index{fair price} $\hat P_t$ is its estimate, at time $t$, of the efficient price $P^\ast_t$ of an instrument
from the information it has at $t$: the books of the venues where the instrument trades, the prices of related instruments, and its own recent fills.
Quotes are set around it; the gap between a fill's price and the fair price at the fill is the fill's expected profit.
\end{definition}

The name is chosen against two neighbours. The \emph{fair value} of a future (One Quant Book 1, chapter 21) is a no-arbitrage price from spot,
financing and dividends; the \emph{theoretical value} of an option (One Quant Book 5, chapter 26) comes from a volatility surface. A \omterm{def:hf:fair-value:fair}{fair price} can
use either as an input, but it is a statistical estimate, revised on every message, and its quality is measured by how well it forecasts.

Three estimators use one venue's top of book. The mid $m_t=\tfrac12(b_t+a_t)$ ignores the queue sizes $Q^b_t$ and $Q^a_t$. The weighted mid
\[
w_t=\frac{a_t\,Q^b_t+b_t\,Q^a_t}{Q^b_t+Q^a_t}
\]
moves towards the ask when the bid queue is the larger, because a large bid queue is less likely to be exhausted first. The microprice (One Quant
Book 7, chapter 8) replaces the fixed formula with a table learned from data: the mid plus the average change of the mid over a horizon, by
bucket of the queue imbalance $I_t=Q^b_t/(Q^b_t+Q^a_t)$. Stoikov showed that the microprice so estimated forecasts short-term prices better than the
mid or the weighted mid.

\begin{example}[Three estimates of one book]\label{ex:hf:fair-value:three}
A book shows 300 shares bid at 99 ticks and 100 offered at 100. The mid is 99.5; the imbalance is 0.75 and the weighted mid 99.75. This chapter's
microprice table, fitted on the first simulated hour with five buckets and a five-second horizon, adds $-0.087$, $-0.007$, $0.022$, $0.091$ and
$0.220$ ticks to the mid for imbalances in $[0,0.2)$, $[0.2,0.4)$, \dots, $[0.8,1]$. Here it adds 0.091 and gives 99.59: the data say the weighted
mid overstates what the imbalance predicts.
\end{example}

\section{Fair price across venues}

When one instrument trades on several venues, each venue's book is a noisy, differently delayed view of one efficient price. A venue whose quote
has not changed for a second tells less than one that changed a millisecond ago; a venue with a thin book is noisier than the primary market.

\begin{definition}[Consolidated fair price]\label{def:hf:fair-value:consolidated}
A \emph{consolidated fair price}\index{consolidated fair price} is a \omterm{def:hf:fair-value:fair}{fair price} that combines the books of every venue on which an instrument
trades, each weighted by how much its prices tell about the efficient price, in view of its liquidity and of how recently it changed.
\end{definition}

A \omterm{def:hf:fair-value:consolidated}{consolidated fair price} is not the national best bid and offer (One Quant Book 1, chapter 9), which is a regulatory construct of protected quotes.
The market maker cares about information, not protection: a venue that rarely trades may carry a stale best price and deserve almost no weight.

The simplest calibration is a regression: the future mid of the instrument, minus its current mid, on the gaps between each source and the mid.
On the chapter's market (a follower B quoted on two venues, B and C) the weights fitted on the first hour are 0.73 on B's own microprice and 0.32 on
C's; they sum to about one, as they should for two measurements of one price, and the second venue earns a third of the weight.

\section{Fair price across instruments}

\begin{definition}[Cross-instrument fair price]\label{def:hf:fair-value:cross}
A \emph{cross-instrument fair price}\index{cross-instrument fair price} is a \omterm{def:hf:fair-value:fair}{fair price} that also uses the prices of related instruments, translated
into the instrument's own units (by a hedge ratio, a conversion ratio, a currency or a fair-value basis), so that a move of a related instrument
that leads is reflected before the instrument's own book shows it.
\end{definition}

An exchange-traded fund's market maker watches the index future; a receipt's market maker watches the home listing and the currency (chapter 14);
an options market maker watches the underlying. The related instrument must be put in the right units first. For the future of an index and a fund
on that index, the translation is the fair-value basis of One Quant Book 1, chapter 21; in the chapter's market the two share one efficient price,
so no translation is needed.

With the leader's mid added to the regression, the weights become 0.74 on B's microprice, 0.22 on C's and 0.18 on the leader's; when the leader
moves two seconds before B, the leader's weight rises to 0.27, B's own to 0.89, and C's falls to $-0.09$. The regression is using the lead that chapter 8 will
estimate and trade.

\section{Filtering: the fair price as a state}

Regressions weigh sources by their average information. A filter weighs each observation when it arrives, by how long it has been since the last
one.

\begin{definition}[Fair-price filter]\label{def:hf:fair-value:filter}
A \emph{fair-price filter}\index{fair-price filter} treats the efficient price as a hidden state, a random walk with variance $q$ per second, and
every source's price as an observation of it with noise variance $r_j$. On an observation $y$ from source $j$ at time $t$, the previous update
having been at $t'$, it propagates and updates:
\[
\Sigma\leftarrow\Sigma+q\,(t-t'),\qquad K=\frac{\Sigma}{\Sigma+r_j},\qquad \hat P_t\leftarrow\hat P_t+K\,(y-\hat P_t),\qquad \Sigma\leftarrow(1-K)\,\Sigma .
\]
\end{definition}

It is the local-level Kalman filter of One Quant Book 4, chapter 19, run on irregular times; a quiet period raises $\Sigma$ and hence the weight of the
next observation, which is exactly the behaviour a stale venue calls for (\cref{lst:hf:fair-value:filter}).

\begin{proposition}[Steady state under regular observations]\label{prop:hf:fair-value:steady}
If one source is observed every $\Delta$ seconds, the posterior variance converges to
\[
\bar\Sigma=\tfrac12\Bigl(-q\Delta+\sqrt{q^2\Delta^2+4q\Delta r}\Bigr),\qquad \bar K=\frac{\bar\Sigma+q\Delta}{\bar\Sigma+q\Delta+r}.
\]
\end{proposition}

\begin{proof}
A fixed point of $\Sigma\mapsto (\Sigma+q\Delta)r/(\Sigma+q\Delta+r)$ satisfies $\Sigma^2+q\Delta\Sigma-q\Delta r=0$; its positive root is
$\bar\Sigma$. The map is increasing
and concave with slope below one at the fixed point, so the iteration converges from any positive start.
\end{proof}

With $q=0.05$ square ticks a second, $r=0.1$ and ten observations a second, $\bar\Sigma=0.02$ and $\bar K=0.2$: each new mid moves the \omterm{def:hf:fair-value:fair}{fair price} by a
fifth of its surprise. The chapter's filter uses $q=0.047$, the variance rate of the efficient price over the first hour, and for each source the
variance of its gap to the mid five seconds later, 0.104, 0.123 and 0.134 for B, C and the leader.

\begin{omfigure}
\begin{tikzpicture}[font=\small,>=stealth,box/.style={draw,rounded corners,align=center,minimum height=1.0cm,text width=2.9cm,fill=omDef!6}]
\node[box] (b) at (0,2.2) {venue B book\\ mid, queues};
\node[box] (c) at (0,0.8) {venue C book\\ mid, queues};
\node[box] (a) at (0,-0.6) {leader A\\ (the index future)};
\node[box,fill=omProp!10,text width=3.4cm] (f) at (5.2,0.8) {regression or filter:\\ weights by information};
\node[box,fill=omThm!8,text width=2.3cm] (p) at (9.6,0.8) {fair price $\hat P_t$\\ of B};
\draw[->] (b.east) -- (f.west);
\draw[->] (c.east) -- (f.west);
\draw[->] (a.east) -- node[below right,font=\scriptsize,pos=0.35] {in B's units} (f.west);
\draw[->] (f) -- (p);
\end{tikzpicture}

\omcaption{The chapter's \omterm{def:hf:fair-value:fair}{fair price} for instrument B: B's own book on its main venue, its book on a second venue, and a leading instrument
translated into B's units, combined by weights that reflect what each tells about B's efficient price.}\label{fig:hf:fair-value:sources}
\end{omfigure}

\section{Evaluating a fair price}

A \omterm{def:hf:fair-value:fair}{fair price} is judged by what it forecasts. In a simulation the efficient price is known and the error can be measured against it; in a live
market only future prices are observed, and the natural target is the mid some seconds later. Both targets are used below, on a grid of 0.1 seconds
over the second simulated hour, with every parameter fitted on the first hour.

\begin{center}
\small
\begin{tabular}{@{}lcccc@{}}
\toprule
& \multicolumn{2}{c}{RMSE (ticks), lead 0.5 s} & \multicolumn{2}{c}{change against the mid (\%)}\\
estimator of B's \omterm{def:hf:fair-value:fair}{fair price} & efficient price & mid in 5 s & lead 0.5 s & lead 2 s\\
\midrule
mid & 0.962 & 0.549 & 0.0 & 0.0\\
weighted mid & 0.963 & 0.542 & $-1.3$ & $+0.1$\\
microprice & 0.957 & 0.536 & $-2.4$ & $-1.7$\\
consolidated, regression & 0.947 & 0.542 & $-1.2$ & $-3.0$\\
consolidated, filter & 0.961 & 0.559 & $+1.8$ & $-0.3$\\
cross-instrument, regression & 0.930 & 0.529 & $-3.5$ & $-12.8$\\
cross-instrument, filter & 0.941 & 0.544 & $-0.8$ & $-8.7$\\
\bottomrule
\end{tabular}
\end{center}

The last two columns are changes of the error against the mid five seconds ahead, averaged over three simulated markets. Three lessons stand out.
The book of the instrument itself says little: in this market every estimator built from it is within 2.5\% of the mid. The leader says more,
and more the further ahead it is (\cref{fig:hf:fair-value:lags}); at a two-second lead the \omterm{def:hf:fair-value:cross}{cross-instrument fair price} removes an eighth of the
error. And the filter, whose variances were set by a rule of thumb, does worse than the regression, which was fitted to the very target it is
scored on: calibrate a \omterm{def:hf:fair-value:fair}{fair price} to what it will be used to forecast.

\begin{omfigure}
\begin{tikzpicture}
\begin{axis}[width=11cm,height=5.2cm,xlabel={\small lead of A over B (seconds)},ylabel={\small change of RMSE (\%)},xmode=log,log basis x=10,
  xmin=0.08,xmax=2.5,ymin=-14,ymax=2,xtick={0.1,0.25,0.5,1,2},xticklabels={0.1,0.25,0.5,1,2},tick label style={font=\footnotesize},
  grid=major,grid style={gray!25},table/col sep=comma,legend style={font=\scriptsize,at={(0.5,-0.32)},anchor=north},legend columns=4]
\addplot[black,thick,mark=*,mark size=1.2pt] table[x=lag,y=micro]{figdata/hft/02-fair-value/lags.csv};
\addplot[omDef,thick,dashed,mark=square*,mark size=1.2pt] table[x=lag,y=consolidated]{figdata/hft/02-fair-value/lags.csv};
\addplot[omThm,thick,mark=triangle*,mark size=1.5pt] table[x=lag,y=cross]{figdata/hft/02-fair-value/lags.csv};
\addplot[omProp,thick,dotted,mark=diamond*,mark size=1.5pt] table[x=lag,y=cross_filter]{figdata/hft/02-fair-value/lags.csv};
\legend{microprice,consolidated,{cross, regression},{cross, filter}}
\end{axis}
\end{tikzpicture}

\omcaption{Change of the error of each \omterm{def:hf:fair-value:fair}{fair price}, against the mid's, as a forecast of B's mid five seconds later, by how many seconds the leader
A moves before B (log scale); mean of three simulated hours each. Data: \texttt{hf\_fairvalue.by\_lag}.}\label{fig:hf:fair-value:lags}
\end{omfigure}

\Cref{fig:hf:fair-value:window} shows why. After the largest ten-second move of the second hour, B's efficient price climbs five ticks, its mid follows
a tick at a time over several seconds, and the \omterm{def:hf:fair-value:cross}{cross-instrument fair price} moves ahead of the mid as soon as the leader's book has moved. In \cref{fig:hf:fair-value:lags}, with three markets per point, the curves are noisy (the half-second lead happens to show less than the quarter-second one); the trend is not.

\begin{omfigure}
\begin{tikzpicture}
\begin{axis}[width=11cm,height=5.2cm,xlabel={\small seconds from the move},ylabel={\small price (ticks)},tick label style={font=\footnotesize},
  xmin=-5,xmax=15,grid=major,grid style={gray!25},table/col sep=comma,scaled y ticks=false,
  yticklabel style={/pgf/number format/.cd,fixed,precision=0,1000 sep={}},
  legend style={font=\scriptsize,at={(0.5,-0.32)},anchor=north},legend columns=3]
\addplot[black,thick,const plot] table[x=t,y=truth]{figdata/hft/02-fair-value/window.csv};
\addplot[omDef,thick,const plot,dashed] table[x=t,y=mid]{figdata/hft/02-fair-value/window.csv};
\addplot[omThm,thick] table[x=t,y=cross]{figdata/hft/02-fair-value/window.csv};
\legend{efficient price,mid,cross-instrument fair price}
\end{axis}
\end{tikzpicture}

\omcaption{Instrument B around the largest ten-second move of its efficient price in the second simulated hour (lead of A over B two seconds): B's efficient price,
its mid and its \omterm{def:hf:fair-value:cross}{cross-instrument fair price} on a 0.1-second grid. Data: \texttt{hf\_fairvalue.window}.}\label{fig:hf:fair-value:window}
\end{omfigure}

\begin{remark}[How good is good]\label{rem:hf:fair-value:good}
A 3\% lower error sounds like nothing. A market maker earns a fraction of the spread on every fill and loses the adverse selection that its
\omterm{def:hf:fair-value:fair}{fair price} failed to foresee (chapter 1); chapter 6 shows the fills that move against it are concentrated in the moments when its \omterm{def:hf:fair-value:fair}{fair price} was
most wrong. The value of a \omterm{def:hf:fair-value:fair}{fair price} is measured in adverse selection saved, per fill, times all the fills of a year.
\end{remark}

\section{Tutorial: three markets, seven fair prices}\label{tut:hf:fair-value:tut}

\begin{tutorial}
\textbf{Goal.} Build the chapter's \omterm{def:hf:fair-value:fair}{fair prices} for a follower instrument quoted on two venues with a leader, fit them on one hour and score them
on the next. \textbf{End state:} the table above and \cref{fig:hf:fair-value:lags,fig:hf:fair-value:window}.
\begin{tutsteps}
\item \textbf{Markets.} \texttt{hf\_fairvalue.markets(lag)} simulates A, and B and C whose efficient price is A's \texttt{lag} seconds later,
each with its own order flow.
\item \textbf{The filter.} \texttt{firm.fairprice.FairFilter} is \cref{def:hf:fair-value:filter} in code; its C++20 and Rust twins reproduce the
same numbers on a shared fixture of 3,000 observations.
\omcode{code/firm/fairprice/firm_fairprice.py}{68}{92}{The fair-price filter: propagate for the elapsed time, then update.}\label{lst:hf:fair-value:filter}
\item \textbf{The regressions.} The future mid minus the mid on the gaps between each source and the mid, fitted on the first hour.
\omcode{code/hft/02-fair-value/python/hf_fairvalue.py}{100}{113}{Regression fair prices: design and weights.}\label{lst:hf:fair-value:weights}
\item \textbf{Score.} \texttt{table(lag)} averages the errors over three seeds; \texttt{fig\_fairvalue.py} writes the figures' data.
\end{tutsteps}
\textbf{What to change next.} Make venue C thin and slow and watch its weight fall; let the leader's lead vary through the day and refit the weights
every ten minutes.
\end{tutorial}

\section{Build: the fair-price engine}\label{bld:hf:fair-value:fairprice}

\begin{build}
\bfield{Purpose} Streaming \omterm{def:hf:fair-value:fair}{fair prices} for every Quoter of the book, with a core fast enough for the hot path.
\bfield{Interface} \texttt{weighted\_mid}, \texttt{imbalance}, \texttt{fit\_micro(imb, spread, dmid, n)}, \texttt{micro(\dots, g)};
\texttt{FairFilter(q, r).update(t, src, y) -> estimate}, \texttt{predict(t)}; \texttt{run\_filter}; \texttt{rmse}.
\bfield{Rules} Parameters are fitted on data before the period they are used on. Sources are translated into the instrument's units before they
enter. The first observation initialises the state. The C++20 (\texttt{cpp/firm\_fairprice.hpp}) and Rust (\texttt{rust/}) filters perform the
same operations in the same order.
\bfield{Acceptance tests} \texttt{code/firm/fairprice/tests/}: the weighted mid and the microprice by hand; one Kalman step by hand; the fixture
reproduced by the Python reference and, in \texttt{cpp/} and \texttt{rust/}, by the twins, to $10^{-9}$ ticks.
\bfield{Stretch} Staleness-dependent noise $r_j(t-t_j)$; a two-state filter with the leader's lead as a state; online refitting of the weights.
\end{build}

\begin{omsources}
\item S.~Stoikov, ``The micro-price: a high-frequency estimator of future prices'', \emph{Quantitative Finance} 18(12), 2018.
\item One Quant Book 4, chapter 19 (the Kalman filter) and One Quant Book 7, chapter 8 (the microprice and the order-book features).
\end{omsources}

\section{Exercises}

\begin{exercise}[$\star$]\label{exo:hf:fair-value:1}
A book shows 300 shares bid at 99 and 100 offered at 100. Compute the weighted mid.
\end{exercise}

\begin{exercise}[$\star$]\label{exo:hf:fair-value:2}
With the chapter's microprice table, what is the microprice of a one-tick book 99/100 whose imbalance is 0.9?
\end{exercise}

\begin{exercise}[$\star$]\label{exo:hf:fair-value:3}
The filter holds $\hat P=100$ with $\Sigma=0.02$; $q=0.05$, $r=0.1$. A mid of 101 arrives 0.4 seconds after the last update. Give the gain and the new
\omterm{def:hf:fair-value:fair}{fair price}.
\end{exercise}

\begin{exercise}[$\star\star$]\label{exo:hf:fair-value:4}
Using \cref{prop:hf:fair-value:steady}, compute $\bar\Sigma$ and $\bar K$ for $q=0.05$, $r=0.1$ and ten observations a second, and say what
happens to $\bar K$ when the source slows to one observation a second.
\end{exercise}

\begin{exercise}[$\star\star$]\label{exo:hf:fair-value:5}
Why can a filter with variances set by rule of thumb forecast worse than the plain mid, while a regression fitted on the same history does not?
\end{exercise}

\begin{exercise}[$\star\star$]\label{exo:hf:fair-value:6}
From \cref{fig:hf:fair-value:lags}, read the improvement of the cross-instrument regression at a one-second lead.
\end{exercise}

\begin{exercise}[$\star\star\star$]\label{exo:hf:fair-value:7}
\emph{Coding.} Compute \texttt{table(1.0)} and give the change of the error of the cross-instrument filter against the mid, five seconds ahead.
\end{exercise}

\begin{exercise}[$\star\star\star$]\label{exo:hf:fair-value:8}
\emph{Find the flaw.} ``Our new microprice table cuts the forecast error by 9\% on today's data, on which we fitted it.''
\end{exercise}

\section{Problem: Whose Price Is It}

\begin{problem}[{Weekend problem --- whose price is it}]\label{pb:hf:fair-value:1}
An exchange-traded fund trades on two venues and follows its index future. You must give its market maker one number, updated on every message.

\medskip\noindent\textbf{Part I --- One book.}
\begin{enumerate}
\item Define the \omterm{def:hf:fair-value:fair}{fair price} and distinguish it from a future's fair value.
\item Write the weighted mid and explain its direction.
\item How is a microprice table fitted, and why only on one-tick books?
\item Give the chapter's table and the microprice of a 99/100 book at imbalance 0.75.
\item By how much does the microprice improve on the mid as a five-second forecast?
\end{enumerate}

\medskip\noindent\textbf{Part II --- Many venues, many instruments.}
\begin{enumerate}[resume]
\item Define a \omterm{def:hf:fair-value:consolidated}{consolidated fair price} and say why it is not the national best bid and offer.
\item Give the regression weights on B's and C's microprices and interpret their sum.
\item Define a \omterm{def:hf:fair-value:cross}{cross-instrument fair price} and say what must happen to the leader's price before it enters.
\item Give the three weights with the leader at a half-second and at a two-second lead.
\end{enumerate}

\medskip\noindent\textbf{Part III --- The filter.}
\begin{enumerate}[resume]
\item Write the \omterm{def:hf:fair-value:filter}{fair-price filter}'s propagate and update steps.
\item Prove \cref{prop:hf:fair-value:steady}.
\item Why does a stale source get more weight on its next observation?
\item Give the chapter's $q$ and $r$ and how each was estimated.
\item Why did the filter forecast worse than the regression?
\end{enumerate}

\medskip\noindent\textbf{Part IV --- The verdict.}
\begin{enumerate}[resume]
\item State the \emph{named result}: the five-second forecast error of the mid, the microprice and the \omterm{def:hf:fair-value:cross}{cross-instrument fair price} at a
half-second lead, and the share of the cross-instrument improvement that the leader brings.
\item How does the improvement change with the lead?
\item Why are three simulated hours too few to trust the shape of \cref{fig:hf:fair-value:lags} point by point?
\item What does \cref{fig:hf:fair-value:window} show?
\item How is a \omterm{def:hf:fair-value:fair}{fair price}'s quality turned into money?
\item In one sentence: what makes a \omterm{def:hf:fair-value:fair}{fair price} better than the mid?
\end{enumerate}
\end{problem}

\section{Interview questions}

\begin{interviewq}[$\star$ \iqroles{trader}]\label{iq:hf:fair-value:1}
The bid queue is ten times the ask queue. Where is the \omterm{def:hf:fair-value:fair}{fair price}, and why?
\end{interviewq}

\begin{interviewq}[$\star\star$ \iqroles{researcher}]\label{iq:hf:fair-value:2}
How would you evaluate a fair-price model when the efficient price is never observed?
\end{interviewq}

\begin{interviewq}[$\star\star$ \iqroles{researcher}]\label{iq:hf:fair-value:3}
Derive the steady-state Kalman gain of a random walk observed with noise every $\Delta$ seconds.
\end{interviewq}

\begin{interviewq}[$\star\star$ \iqroles{developer}]\label{iq:hf:fair-value:4}
A fair-price service consumes feeds from eight venues. What must each update carry, and what do you do when one feed goes silent?
\end{interviewq}

\begin{interviewq}[$\star\star$ \iqroles{trader}]\label{iq:hf:fair-value:5}
Your fund's \omterm{def:hf:fair-value:fair}{fair price} uses the index future. The future's market closes early today. What changes?
\end{interviewq}

\begin{interviewq}[$\star\star\star$ \iqroles{researcher}]\label{iq:hf:fair-value:6}
Two venues' mids measure one price with independent noises of variance $r_1$ and $r_2$. What weights minimise the variance of the combination,
and when is ignoring the second venue optimal?
\end{interviewq}
